\textbf{Chi tiết cài đặt}

Thí sinh cần cài đặt hàm sau:

\begin{lstlisting}
string solve()
\end{lstlisting}

\begin{itemize}
   \item Hàm này có thể được gọi nhiều lần bởi trình chấm.
   \item Hàm này cần trả về xâu $U$ thỏa mãn điều kiện đề bài. Lưu ý là bạn không được biết trước giá trị của $N,L,R$.
\end{itemize}

Trong hàm \texttt{solve()}, thí sinh được phép gọi hàm sau (thí sinh cần khai báo trước thư viện \texttt{smatchlib.h} để có thể sử dụng hàm này):

\begin{lstlisting}
int ask(string s)
\end{lstlisting}

\begin{itemize}
   \item $s$: Xâu mà Bob sẽ đưa cho Alice để thực hiện câu hỏi. Xâu $s$ chỉ được chứa các ký tự Latin in thường và có độ dài không quá $200$.
   \item Hàm sẽ trả về một số nguyên không âm là số lượng xâu $t\in \Omega$ thỏa mãn $s$ khớp với $t$ chia lấy dư cho $379$.
\end{itemize}

Ngoài ra, thí sinh được phép cài đặt các biến, các hàm toàn cục, nhưng không được phép có hàm \texttt{main()}.

Cách cài đặt và trình chấm ví dụ có thể xem tại contest material, lưu ý trình chấm ví dụ chỉ mô phỏng cách hoạt động và sẽ không được dùng để chấm chính thức.

\textbf{Ràng buộc}

\begin{itemize}
   \item $4\le N\le 13$
   \item $1\le L\le R\le N$
\end{itemize}

Trình chấm của bài này \textbf{không mang tính thích ứng}. Điều này tương đương với việc các giá trị $N,L,R$ sẽ được cố định từ đầu và không thay đổi trong suốt quá trình chương trình của bạn tương tác với trình chấm.